\originalchapter{对称性与守恒量}
\originalsection{循环坐标与能量}
\begin{definition}
广义动量是$p_i=\partial L/\partial\dot q_i$. 若$L$不显含$q_i$, 则$q_i$为循环坐标. 由Euler--Lagrange方程有$\dot p_i=0$.
\end{definition}
广义动量不一定是$m\dot q_i$. 例如平面极坐标$T=m(\dot r^2+r^2\dot\phi^2)/2$给$p_\phi=mr^2\dot\phi$, 它是角动量, 单位也与直线动量不同. 判断守恒要看拉格朗日量的坐标依赖, 不能只看“速度是否不变”.

定义能量函数$E=\sum_ip_i\dot q_i-L$. 沿运动方程计算,
\[
\dot E=\sum_i(\dot p_i-L_{q_i})\dot q_i-L_t=-L_t.
\]
所以$L$不显含时间时$E$守恒. 当约束固定, $T$是速度的齐次二次式, 势能与速度无关, Euler齐次公式给$\sum_iT_{\dot q_i}\dot q_i=2T$, 从而$E=T+U$. 如果势依赖速度或者坐标系随时间移动, 后一个等式可能失效, 前面的能量函数推导却仍成立. 例如$U(q,\dot q)$不显含$t$时, $E=T+U-\sum_i\dot q_iU_{\dot q_i}$.

\originalsection{Noether定理}
对称性是“改变描述或位置后作用量不变”的精确陈述. 不仅严格不变有用, 变一个端点项也足以产生守恒量.
\begin{theorem}
设一族无穷小变换为$\delta t=\epsilon\tau(q,t)$, $\delta q_i=\epsilon\xi_i(q,t)$. 若拉格朗日密度的一阶变化满足
\[
\sum_i(L_{q_i}\xi_i+L_{\dot q_i}(\dot\xi_i-\dot q_i\dot\tau))+L_t\tau+L\dot\tau=\dot B,
\]
则沿Euler--Lagrange方程, $J=\sum_ip_i\xi_i-E\tau-B$守恒.
\end{theorem}
\begin{proof}
将$\dot\xi_i-\dot q_i\dot\tau$代入左端, 再用$L_{q_i}=\dot p_i$, $\dot E=-L_t$和$\sum_ip_i\dot q_i=E+L$, 左端成为
$\dd(\sum_ip_i\xi_i-E\tau)/\dd t$. 与$\dot B$相等, 故$\dot J=0$.
\end{proof}
若平移不改变$L$, 取$\xi_a=\vect e$, $\tau=B=0$, 得总动量在$\vect e$方向的分量守恒. 若转动不改变$L$, 取$\delta\vect r_a=\epsilon\vect e\times\vect r_a$, 得$J=\vect e\cdot\sum_a\vect r_a\times\vect p_a$, 即角动量. 若时间平移不改变$L$, 取$\tau=1$, $\xi=B=0$, 得$J=-E$. 负号不影响守恒. 这些例子说明能量守恒对应时间均匀性, 动量对应空间平移, 角动量对应旋转.

\originalsection{两体问题与质心}
对于$L=\frac12m_1\dot{\vect r}_1^2+\frac12m_2\dot{\vect r}_2^2-U(|\vect r_1-\vect r_2|)$, 定义$M=m_1+m_2$, $\mu=m_1m_2/M$, $\vect R=(m_1\vect r_1+m_2\vect r_2)/M$, $\vect r=\vect r_1-\vect r_2$. 逆变换为$\vect r_1=\vect R+(m_2/M)\vect r$, $\vect r_2=\vect R-(m_1/M)\vect r$. 代入动能时交叉项系数$m_1m_2/M-m_2m_1/M$为零, 所以
\[
L=\frac12M\dot{\vect R}^2+\frac12\mu\dot{\vect r}^2-U(r).
\]
质心作匀速直线运动, 相对运动等价于质量$\mu$的单粒子中心力运动. 只有在重粒子质量远大于另一粒子时, 才可进一步用轻粒子的质量近似$\mu$并将重粒子固定.

\originalsection{位力定理}
令$G=\sum_a\vect r_a\cdot\vect p_a$, 对通常的$T=\sum_ap_a^2/(2m_a)$有$\dot G=2T+\sum_a\vect r_a\cdot\vect F_a$. 若运动有界且$G$有界, 则长时平均$\langle\dot G\rangle=\lim_{\mathcal T\to\infty}(G(\mathcal T)-G(0))/\mathcal T=0$. 对齐次势$U(\lambda\vect r)=\lambda^kU(\vect r)$, 微分得$\sum_a\vect r_a\cdot\nabla_aU=kU$, 因而$2\langle T\rangle=k\langle U\rangle$. 周期运动可直接作一个周期平均. 谐振子$k=2$给$\langle T\rangle=\langle U\rangle$; 引力势$k=-1$给$2\langle T\rangle=-\langle U\rangle$. 无界散射轨道不能不加条件就套这个平均公式.

本章来源: 8.223第5--6讲、8.09第1章及两课相关作业. Noether定理的一般时间变换版本和完整平均条件为中文补充.
